\section{System Design and Architecture}
\label{sec:system-design}

\subsection{Configuration-Centric Architecture}
\label{sec:config-centric-architecture}

The core design choice in this repository is that experimentation is \emph{schema first}. Instead of exposing a large number of loosely checked flags, the project forces model topology, optimizer family, training limits, and telemetry options through a single validated configuration document. This reduces ambiguity when reproducing results and makes it possible to compare many architectures under a consistent operational interface.

The authoritative contract is \texttt{configs/schema.yaml}. It enforces three top-level objects:
\begin{itemize}[leftmargin=*]
    \item \texttt{model\_class}
    \item \texttt{model}
    \item \texttt{training}
\end{itemize}

\paragraph{Model Class Selection.}
The \texttt{model\_class} field determines the architectural variant instantiated by the training pipeline. Two options are supported:
\begin{itemize}[leftmargin=*]
    \item \textbf{frankenstein}: Mixed-architecture encoder models supporting diverse attention mechanisms (standard, sigmoid, retentive, state-space, sparse, and gated mixers) with MoE (Mixture of Experts) and advanced features. Optimized for bidirectional encoder-style training with masked language modeling (MLM) objectives.
    \item \textbf{frankensteindecoder}: Autoregressive causal decoder for LLM-style next-token generation. Enables causal attention masking for sequential text generation tasks. When this class is selected, runtime enforces \texttt{mode='decoder'}. 
\end{itemize}

\paragraph{Training Mode Selection.}
The \texttt{model.mode} field controls attention masking behavior across the model:
\begin{itemize}[leftmargin=*]
    \item \textbf{encoder}: Uses bidirectional attention where all tokens attend to all other tokens in the sequence. Suitable for masked language modeling (MLM) pre-training tasks where the model learns to predict randomly masked tokens based on full context.
    \item \textbf{decoder}: Uses causal masking where each token can only attend to previous tokens in the sequence. Required for autoregressive (AR) next-token prediction tasks such as language modeling and text generation. When \texttt{model\_class='frankensteindecoder'}, the system automatically forces \texttt{mode='decoder'} at runtime.
\end{itemize}

This dual architecture support enables the system to handle both encoder-style pre-training (MLM on bidirectional contexts) and decoder-style generation (causal autoregressive prediction) through a unified configuration interface.

\textbf{Sequence Mixer Pattern Selection}: The \texttt{layer\_pattern} field accepts an ordered list of mixer identifiers spanning seven families: dense baselines (\texttt{standard\_attn}, \texttt{sigmoid\_attn}, \texttt{gqa}), recurrent/retentive (\texttt{retnet}, \texttt{retnet\_attn}, \texttt{mamba}, \texttt{ode}, \texttt{titan\_attn}), sparse attention (\texttt{sparse\_transformer\_attn}, \texttt{longformer\_attn}, \texttt{bigbird\_attn}, \texttt{sparsek\_attn}, \texttt{nsa\_attn}, \texttt{sparge\_attn}, \texttt{fasa\_attn}, \texttt{msa\_attn}, \texttt{sparda\_attn}), gated mechanisms (\texttt{gla\_attn}, \texttt{deltanet\_attn}, \texttt{gated\_deltanet\_attn}, \texttt{gated\_deltanet2\_attn}, \texttt{hgrn2\_attn}, \texttt{fox\_attn}, \texttt{gated\_softmax\_attn}, \texttt{kda\_attn}), latent compression (\texttt{mla\_attn}, \texttt{gqla\_attn}, \texttt{mlra\_attn}, \texttt{tucker\_attn}, \texttt{iha\_attn}, \texttt{gta\_attn}, \texttt{mtla\_attn}, \texttt{cca\_attn}, \texttt{ccgqa\_attn}, \texttt{gma\_attn}), conditional memory (\texttt{engram\_attn}), and fast-weight attention (\texttt{falcon1\_attn}, \texttt{falcon2\_attn}, \texttt{falcon3\_attn}, \texttt{falcon1a\_attn}, \texttt{falcon2a\_attn}, \texttt{falcon3a\_attn}). The complete taxonomy with references is provided in Section~\ref{sec:architecture-taxonomy} and Appendix~\ref{sec:annex-summary-tables}.

The \texttt{training.optimizer.optimizer\_class} supports a broad optimizer family:
\texttt{sgd\_momentum}, \texttt{adamw}, \texttt{adafactor}, \texttt{galore\_adamw}, \texttt{prodigy}, \texttt{lion}, \texttt{sophia}, \texttt{muon}, \texttt{turbo\_muon}, \texttt{radam}, \texttt{adan}, \texttt{adopt}, \texttt{ademamix}, \texttt{mars\_adamw}, \texttt{cautious\_adamw}, \texttt{lamb}, \texttt{schedulefree\_adamw}, \texttt{shampoo}, \texttt{soap}, \texttt{anon}, \texttt{apollo}, \texttt{apollo\_mini}, and \texttt{q\_apollo}.

\begin{figure}[t]
\centering
\fitfigure{\begin{tikzpicture}[
font=\scriptsize,
box/.append style={inner sep=4pt, minimum height=0.62cm}
]
% Top level
\node[box, fill=green!10, minimum width=1.9cm] (yaml) at (0,0) {\textbf{YAML Config}};
\node[box, fill=green!8, minimum width=2.0cm] (validation) at (0,-1.55) {\textbf{Validation}\\\texttt{schema + rules}};
\node[box, fill=blue!10, minimum width=1.9cm] (web) at (-3.9,-2.95) {\textbf{Web}\\\texttt{Streamlit}};
\node[box, fill=blue!8, minimum width=2.1cm] (cli) at (3.9,-2.95) {\textbf{CLI}\\\texttt{train / deploy / quantize / infer / sbert-train / sbert-infer / transformers-export / bitnet-gguf / web-server}};
% Main blocks
\node[box, fill=orange!12, minimum width=1.9cm] (model) at (-2.6,-4.3) {\textbf{Model}\\35+ mixers};
\node[box, fill=yellow!15, minimum width=2.0cm] (train) at (0,-4.3) {\textbf{Training}\\AMP + scheduler};
\node[box, fill=red!12, minimum width=1.9cm] (opt) at (2.6,-4.3) {\textbf{Optimizer}\\23 families};
% Outcomes
\node[box, fill=purple!12, minimum width=2.0cm] (modelout) at (-2.6,-5.95) {\textbf{Model Build}\\pattern + loops};
\node[box, fill=purple!10, minimum width=2.0cm] (trainout) at (0,-5.95) {\textbf{Train Loop}\\checks + logs};
\node[box, fill=purple!10, minimum width=2.0cm] (optout) at (2.6,-5.95) {\textbf{Opt Step}\\group routing};
% Deployment
\node[box, fill=cyan!12, minimum width=2.0cm] (deploy) at (0,-7.75) {\textbf{Deploy}\\quantization};
\node[box, fill=cyan!10, minimum width=2.0cm] (sbert) at (2.6,-7.75) {\textbf{SBERT}\\search/cluster};
\node[box, fill=cyan!8, minimum width=2.0cm] (export) at (-2.6,-7.75) {\textbf{Export}\\transformers-export / bitnet-gguf};
% Arrows
\draw[line] (yaml) -- (validation);
\draw[line] (validation) -- (web);
\draw[line] (validation) -- (cli);
\draw[line] (validation) -- (model);
\draw[line] (validation) -- (train);
\draw[line] (validation) -- (opt);
\draw[line] (model) -- (modelout);
\draw[line] (train) -- (trainout);
\draw[line] (opt) -- (optout);
\draw[line] (trainout) -- (deploy);
\draw[line] (optout) -- (sbert);
\draw[line, dashed] (modelout) -- (trainout);
\draw[line, dashed] (optout.south) |- (deploy.east);
\draw[line] (cli) -- (export);
\end{tikzpicture}}
\caption{System architecture: configuration flows through validation, partitions into model/training/optimizer, executes runtime, produces deployment/SBERT artifacts.}
\label{fig:system-architecture}
\end{figure}

\subsection{Schema Scope and Validation Rules}
\label{sec:schema-validation}

The schema is strict: top-level and nested objects set \texttt{additionalProperties: false}. This guarantees that unknown keys fail fast instead of being silently ignored. The \texttt{training.optimizer.parameters} object is additionally constrained by optimizer-specific prefix rules through \texttt{allOf}+\texttt{if/then} pattern checks.

Normalization values currently accepted by schema are:
\[
\texttt{norm\_type} \in \{\texttt{layer\_norm}, \texttt{dynamic\_tanh}, \texttt{derf}, \texttt{rms\_norm}, \texttt{prms\_norm}\}
\]

\subsection{Configuration Schema}
The complete model and training configuration schema, including all fields, types, ranges, and validation rules, is documented in Appendix~\ref{sec:annex-schema}. The schema enforces \texttt{additionalProperties: false} at all levels, guaranteeing that unknown keys fail fast.

Looped depth induced by schema is:
\[
L_{\text{logical}} = \texttt{num\_layers} \times \texttt{num\_loops}
\]
which is the configuration-level definition of looped blocks.

\subsection{Training Task Types}

The \texttt{training.task} field determines the training objective, working in conjunction with \texttt{model.mode} to define how the model learns:

\paragraph{Masked Language Modeling (MLM).}
\begin{itemize}[leftmargin=*]
    \item \textbf{Mode}: Requires \texttt{mode='encoder'} for bidirectional attention.
    \item \textbf{Objective}: Randomly mask tokens in input sequence (typically 15\%) and train model to predict masked tokens based on full bidirectional context.
    \item \textbf{Use Case}: Pre-training encoders for representation learning, following BERT-style methodology \citep{devlin_bert_2019}. Model learns bidirectional representations capturing context from both left and right.
    \item \textbf{Configuration}: Uses \texttt{mlm\_probability} parameter to control masking fraction.
\end{itemize}

\paragraph{Causal Language Modeling (\texttt{causal\_lm}).}
\begin{itemize}[leftmargin=*]
    \item \textbf{Mode}: Requires \texttt{mode='decoder'} for causal masking.
    \item \textbf{Objective}: Autoregressive next-token prediction --- each position $i$ predicts token $i+1$ given all preceding tokens only. Cross-entropy is computed over all shifted non-pad positions (no masking), in contrast to MLM which supervises only randomly masked positions.
    \item \textbf{Use Case}: Language generation and LLM-style tasks following GPT methodology. Model learns sequential dependencies with causal attention where each token can only attend to preceding tokens.
    \item \textbf{Model Class}: Set \texttt{task='causal\_lm'} and \texttt{model\_class='frankensteindecoder'}; the schema and runtime require the decoder class for this task (causal masking is only provided by the decoder).
    \item \textbf{Configuration}: Reuses \texttt{training.optimizer}; the streaming dataset stores sequences unmasked (\texttt{labels == input\_ids}) so the trainer can compute the shifted next-token loss.
\end{itemize}

This task support enables unified experimentation across encoder-style pre-training (MLM for bidirectional understanding), decoder-style generation (\texttt{causal\_lm} for autoregressive next-token prediction), and sentence embedding (SBERT), all within the same codebase.


\subsection{Optimizer Configuration}
The \texttt{training.optimizer} object selects among 23 optimizer families via \texttt{optimizer\_class} and provides per-parameter-group hyperparameter control through a prefixed naming convention. The complete optimizer prefix contract and supported families are documented in Appendix~\ref{sec:annex-schema}.

\subsection{Training Safety and Runtime Semantics}

Schema-level safety features include accumulation, clipping, post-clip explosion checks, and NaN retries:
\[
g_{\text{acc}}=\frac{1}{K}\sum_{i=1}^{K} g_i,\quad K=\texttt{gradient\_accumulation\_steps}
\]
\[
g_{\text{clip}} = g_{\text{acc}}\cdot
\min\left(1,\frac{\tau}{\|g_{\text{acc}}\|_2+\epsilon}\right),\quad
\tau=\texttt{grad\_clip\_max\_norm}
\]
then overflow guards use \texttt{inf\_post\_clip\_threshold} and retry logic bounded by \texttt{max\_nan\_retries}.

The complete pseudocode for the schema-driven training step---including gradient accumulation, global-norm clipping, post-clip explosion guards, bounded NaN/Inf retries, scheduler dispatch, and rolling/best checkpoint rotation---is deferred to Appendix~\ref{sec:annex-norm-bitnet-mod}.

\subsection{Normalization Variants}
\label{sec:normalization}

Normalization controls activation scale across depth and is also a schema compatibility question, since the accepted \texttt{norm\_type} values are \texttt{layer\_norm}, \texttt{dynamic\_tanh}, \texttt{derf}, \texttt{rms\_norm}, and \texttt{prms\_norm} (partial RMSNorm). The mathematical formulations of LayerNorm, RMSNorm, pRMSNorm, Dynamic Tanh (DyT) \citep{zhu_transformers_2025}, and Dynamic Erf (Derf) \citep{chen_stronger_2025}, together with a comparison table, are detailed in Appendix~\ref{sec:annex-norm-bitnet-mod}.

\subsection{Residual-Stream Architectures}
\label{sec:residual-stream}

The residual connection is not only a modeling choice but also a macro-design decision. The system exposes four residual-connection strategies via the \texttt{model.residuals} sub-object (\texttt{residual\_type} flat key, default \texttt{standard}):

\begin{itemize}[leftmargin=*]
    \item \textbf{Standard residual} (\texttt{residual\_type="standard"}): the classical identity residual \citep{he_deep_2016}, $x_{l+1} = x_l + F(x_l, W_l)$. Stateless, no extra parameters.
    \item \textbf{No residual} (\texttt{residual\_type="none"}, experimental): $x_{l+1} = F(x_l, W_l)$, dropping the skip connection entirely. Useful only as an ablation probe.
    \item \textbf{Attention Residuals} (\texttt{residual\_type="full\_attn"} or \texttt{"block\_attn"}) \citep{sun_attnres_2026}: replace the fixed unit-coefficient residual sum with learned softmax attention over depth. Each layer attends over all previous layer outputs (Full AttnRes) or over a small set of block-wise accumulated representations (Block AttnRes). Adds $L \cdot C$ parameters (one learned query vector $\mathbf{w}_l$ per layer). Zero-init queries make the initial attention an equal-weight average, matching the standard residual at step zero.
    \item \textbf{Manifold-Constrained Hyper-Connections} (\texttt{residual\_type} combined with \texttt{mhc.enabled=true}, documented separately in Appendix~\ref{sec:annex-mhc}): the n-stream residual \citep{xie_mhc_2025} whose stream-mixing matrix is constrained to the Birkhoff polytope via Sinkhorn--Knopp.
\end{itemize}

AttnRes integrates orthogonally with mHC and with Mixture-of-Depths: the \texttt{model.residuals.mhc\_stream\_mode} field selects whether the depth-wise attention runs per stream (\texttt{independent}) or jointly over the flattened $nC$ projection (\texttt{joint}) when mHC is active. The full formulation, the four strategies, the implementation wiring, and the reported results are detailed in Appendix~\ref{sec:annex-attnres}. mHC alone is detailed in Appendix~\ref{sec:annex-mhc}.

